\documentclass[11pt]{article}
\usepackage[T1]{fontenc}
\usepackage[utf8]{inputenc}
\usepackage[english]{babel}
\usepackage[margin=1in]{geometry}
\usepackage{amsmath,amssymb}
\usepackage{enumitem}
\usepackage{microtype}
\usepackage[hidelinks]{hyperref}

\setlength{\parindent}{0pt}
\setlength{\parskip}{0.65em}
\setlist{leftmargin=*,topsep=0.3em,itemsep=0.35em}

\title{Technical Review of\\
\emph{Early-Stopping Activation Steering for Factual Preference Alignment in Vietnamese Medical Language Models}}
\author{}
\date{September 3, 2026}

\begin{document}
\maketitle

\section*{Overall assessment}

The revised manuscript adds useful implementation detail and a paired table for the Hybrid RRF comparison.  The strongest internally supported result remains the Hard Cutoff improvement in the natural-completion evaluation: the displayed discordant counts yield the reported exact McNemar result.  The new Hybrid RAG contingency table is also arithmetically coherent.

The revision nevertheless introduces several definite errors and unresolved scientific inconsistencies.  Most importantly, the abstract, contribution list, and Results report incompatible activation-norm experiments; the covariance-control sample size changes from 20 to 100 without a corresponding method or seed update; the new implementation paragraph contradicts itself about precision and layer numbering; and the LaTeX source does not compile in halt-on-error mode.  A new ``Experiment 10'' claim has no experiment or table in the paper and contains an approximately eightfold error in the stated reduction factor.  The paper still needs substantial technical revision before its mechanistic and comparative conclusions are reliable.

\section*{Major technical findings}

\subsection*{1. The manuscript reports two incompatible activation-norm experiments}

\textbf{Classification: definite internal inconsistency.}

\textbf{Location.} Abstract; contribution 3; Empirical Activation-Norm Dynamics; Discussion \& Limitations.

\textbf{Problem.} The abstract and Results retain the earlier free-generation values: Linear Decay reaches 51.20 at $t=50$, the baseline is 54.21, and Continuous Steering is 56.18.  Contribution 3 instead claims a fixed teacher-forcing experiment in which the baseline is 53.39, Continuous Steering is 56.34, and early stopping returns exactly to 53.39.  The Results section contains no teacher-forcing protocol, table, uncertainty, or these revised values.  Thus readers cannot determine which experiment supports the paper's current mechanistic claim.

Under fixed teacher forcing, equality between the baseline and an early-stopping hook at $t=50>K$ is also largely expected: $\alpha(50)=0$ at the measured layer and the input tokens are held fixed.  Continuous Steering's larger post-hook norm similarly follows directly from adding a nonzero vector.  These observations alone do not show that continuous steering produces harmful cumulative inflation or that early stopping improves generated behavior.

\textbf{Why it matters.} Activation-norm mitigation is the stated motivation and a named contribution.  Conflicting designs and values prevent verification, while a post-hook offset under teacher forcing does not establish the claimed causal link to repetition, syntax, or factual preference.

\textbf{Suggested fix.} Replace all three descriptions with one clearly identified experiment.  State whether decoding is free-running or teacher-forced, which token sequence is used, and how prompts are selected.  Report pre-hook and post-hook norms, $\|\hat h-h\|_2$, $\langle h,v_{\text{steer}}\rangle$, sample counts, and uncertainty over the full trajectory.  Separate the tautological observation that $\alpha(t)=0$ after $K$ from evidence that early stopping improves output quality.

\subsection*{2. The new implementation protocol contains precision, layer-index, and token-timing contradictions}

\textbf{Classification: definite errors and implementation ambiguity.}

\textbf{Location.} Experimental Setup \& Reproducibility Implementation Protocol; Equations~(1), (3), and the added hook equation; abstract and contribution 3.

\textbf{Problem.} The first sentence says NF4 uses \texttt{bfloat16} compute, while the Reproducibility Environment immediately specifies \texttt{float16}.  Both cannot describe the same runs.  The hook target is written as \texttt{model.model.layers[8]} and called the ``0-indexed 8th layer''; index 8 is the ninth decoder block under zero-based indexing.  Elsewhere the intervention is simply called Layer~8, so the scientific layer label is ambiguous by one block.

There is also an autoregressive timing issue.  If prompt prefill is unsteered, the first output token is sampled from unsteered prefill logits.  A hook applied when the first generated token is fed back into the cached model affects the logits for the second output token.  Therefore a counter with $t=1$ on the first post-prefill call does not literally steer output tokens 1 through $K$; it ordinarily affects the predictions for tokens 2 through $K+1$.  The manuscript's statement that steering is restricted to the ``first $K$ generated tokens'' is not consistent with the described prefill behavior unless a different counter/logit convention is implemented.

Equation~(1) still defines a generic \texttt{LayerOutput}, while the new paragraph identifies the post-block tuple element \texttt{layer\_output[0]}.  It remains unclear whether extraction and injection use exactly the same tensor location.

\textbf{Why it matters.} Compute dtype, decoder block, and generation-step alignment can change every reported output.  The current description is not sufficiently self-consistent to reproduce the intervention.

\textbf{Suggested fix.} Choose and report the actual compute dtype from the run configuration.  Write either ``decoder block at zero-based index 8 (the ninth block)'' or change the code/index to match Layer~8.  Provide short executable pseudocode showing the prefill call, when the counter increments, which logits are affected, how the tuple is reconstructed, and whether extraction uses the same post-block hidden state.  Define Continuous and Hard Cutoff piecewise as well as Linear Decay.

\subsection*{3. Covariance-control metadata are mutually inconsistent}

\textbf{Classification: definite reporting inconsistency.}

\textbf{Location.} Directional \& Manifold Covariance Controls; Table~V caption and Cov-Matched row; directional-specificity paragraph.

\textbf{Problem.} The method specifies $N=20$ covariance-matched vectors with seeds 3000--3019 and 10,000 test passes, which equals $20\times500$.  Table~V now reports $N=100$ covariance-matched vectors but retains the same mean and sample SD, 74.81\%$\pm$0.95\%.  If 100 vectors were actually evaluated, the stated seed range and number of passes must change to 100 seeds and 50,000 prompt--vector evaluations.

The prose then labels the learned vector ``directionally superior'' to the covariance-matched distribution because it is 2.39 points, or 2.52 sample SDs, above its mean.  A standardized distance from the mean is descriptive; it is not an empirical tail probability, and the distribution's range or exceedance count is not given.  All vectors also share the same 500 questions, so the statistical unit and dependence structure need to be stated.

\textbf{Why it matters.} The covariance controls are the strongest evidence against a generic on-manifold perturbation explanation.  The current metadata do not establish whether the experiment was expanded or only relabeled.

\textbf{Suggested fix.} Reconcile the actual number of vectors, seed range, evaluation count, mean, SD, and BERTScore aggregation from the run artifacts.  Report the full vector-level distribution and the number of controls exceeding $+v_{\text{steer}}$.  Use a pre-specified Monte Carlo comparison of paired improvements over baseline rather than treating 2.52 SDs as proof of specificity.

\subsection*{4. The pair-shuffled control remains mathematically invalid and is no longer reported}

\textbf{Classification: definite control-design error.}

\textbf{Location.} Contrastive Vector Estimation \& Directional Controls; Table~V.

\textbf{Problem.} The method still lists Pair-Shuffled Steering and states that its cosine with $v_{\text{steer}}$ is 1.0000 because of commutativity.  For the mean-difference estimator,
\[
 \frac{1}{N}\sum_i\left(h_i^+-h_{\pi(i)}^-\right)
 =\mu_l^+-\mu_l^-,
\]
so pairing permutations cannot destroy the learned direction.  The revised Table~V removes the Pair-Shuffled row entirely, although the method continues to count it as a control.  The section also says that ``two controls'' are evaluated, but its second list item contains five conditions and repeats Negative Steering.

The Label-Shuffled control remains ambiguous: flipping only $N_{\text{flipped}}=185$ of the stated 10,290 training pairs would leave most labels unchanged, whereas using a 185-pair construction subset is not documented.

\textbf{Why it matters.} An invariant operation is not a negative control.  Listing an unreported control and an unclear partial-label shuffle overstates the directional validation.

\textbf{Suggested fix.} Delete Pair-Shuffled Steering or replace it with a full balanced permutation of the positive/negative labels or random signs on pairwise differences.  State the complete construction population, number flipped, seeds, and selection rule.  Use a flat list in which every named control has a corresponding result.

\subsection*{5. The new paired Hybrid RAG table does not support an ``outperforms'' conclusion}

\textbf{Classification: unsupported comparative claim.}

\textbf{Location.} Table~VI (Linear Decay versus Hybrid RRF); RAG Results; abstract and conclusion.

\textbf{Problem.} The new contingency table has 27 questions preferred only under steering and 22 preferred only under Hybrid RAG.  The two-sided exact McNemar test is
\[
 p=0.5682,
\]
so the five-question, 1.00-point difference is not statistically established.  In addition, Hybrid RAG has higher absolute reference BERTScore F1 (0.7190 versus 0.7150).  Steering has lower latency and context use, but it does not dominate Hybrid RAG across the reported metrics.

The manuscript continues to say that steering ``outperforms Hybrid RAG'' and emphasizes the weaker BM25 comparator in the abstract and conclusion.  Table~VI is not cited or interpreted in the prose.

\textbf{Why it matters.} The added paired data directly show that the accuracy difference may be sampling variation.  A generic superiority claim is inconsistent with both the test and the BERTScore column.

\textbf{Suggested fix.} Report the exact McNemar result and a paired interval, and say that steering and Hybrid RAG have statistically indistinguishable RefPref rates on this test while exhibiting a quality--latency--context trade-off.  Make Hybrid RAG the principal non-oracle comparator; retain BM25 as a weaker lexical baseline.

\subsection*{6. ``Experiment 10'' is absent, and its $11\times$ claim is arithmetically wrong}

\textbf{Classification: definite arithmetic error and unsupported result.}

\textbf{Location.} Multi-Metric Pareto Trade-off \& Deployment Mitigation.

\textbf{Problem.} The Discussion introduces ``Experiment 10'' with a repetition penalty $\theta_{\text{rep}}=1.15$, Rep-4 of 3.76\%, and latency values of 78.74~s and 45.05~s.  No Experiment 10, method, table, test set, schedule, output length, RefPref result, or uncertainty appears elsewhere in the manuscript.  The paragraph then repeats the original unpenalized discussion almost verbatim.

The statement that Rep-4 ``drops by over $11\times$'' from 5.35\% to 3.76\% is false:
\[
 \frac{5.35}{3.76}=1.42,
 \qquad
 \frac{5.35-3.76}{5.35}=29.7\%.
\]
The stated latency reduction is approximately correct: $(78.74-45.05)/78.74=42.8\%$.  However, lower latency may simply reflect shorter generations induced by the penalty, and the paper does not report EOS or token counts.  It also does not show that the $+5.20$-point RefPref gain is preserved.

\textbf{Why it matters.} This new result is used to claim that the trade-off is ``resolved completely'', but it is neither documented nor arithmetically described correctly.

\textbf{Suggested fix.} Remove the claim until the experiment is added as a complete Results subsection.  Report all four natural-completion conditions with and without the penalty, including RefPref, reference BERTScore, Rep-4, EOS, generated-token count, and synchronized latency.  Replace ``over $11\times$'' with the correct 29.7\% relative reduction if those values are verified.

\subsection*{7. The mechanistic narrative remains inconsistent with the reported repetition results}

\textbf{Classification: unsupported causal claim.}

\textbf{Location.} Introduction; Empirical Activation-Norm Dynamics; natural-completion results; Discussion.

\textbf{Problem.} The Introduction says a persistent norm offset is ``manifested as repetitive text loops and degraded syntax'', but no syntax metric or loop examples are reported.  Under unpenalized natural completion, Continuous Steering has Rep-4 of 4.38\%, whereas Hard Cutoff and Linear Decay have 5.35\% and 5.37\%.  Thus the only reported repetition metric is worse for the early-stopping schedules.  The revised Discussion acknowledges this observation but then speculates that steering focuses attention on ``precise clinical templates'' without an analysis supporting that mechanism.

\textbf{Why it matters.} The empirical results may still support a RefPref benefit, but they do not support the claimed chain from continuous norm offset to repetition and from early stopping to mitigation.

\textbf{Suggested fix.} Separate the empirical RefPref result from the mechanistic hypothesis.  Add matched repetition, syntax, degeneration, and token-diversity analyses before claiming a causal mechanism.  Otherwise state only that continuous intervention produces a post-hook norm offset and that the relationship between this offset and output quality remains unresolved.

\subsection*{8. Benchmark validity and the RefPref endpoint still limit factual and clinical claims}

\textbf{Classification: likely validity issue; parts need external verification.}

\textbf{Location.} Benchmark Construction; BERTScore Reference-Preference Criterion; abstract, Results, and conclusion.

\textbf{Problem.} The paper does not document the qualifications of the claimed experts, question-construction procedure, synthetic-negative generator and prompts, factual review, agreement, or adjudication.  The category counts 165, 160, and 175 sum to 500 and therefore appear to describe the evaluation subset, but they are introduced as if they characterize all 14,700 records.  The relationship between the 2,205-record test partition and the ``question-disjoint'' 500-question subset remains unclear, as does overlap by drug, entity, or source passage.

RefPref measures whether an output is semantically closer to one reference than to one synthetic counter-answer.  It does not verify dosage arithmetic, contraindications, interactions, or new unsupported claims.  Tie counts and complete BERTScore settings are not reported.  The paper acknowledges that RefPref is a proxy but continues to use ``factual accuracy'', ``factual alignment'', and deployment-oriented language.

\textbf{Why it matters.} Source overlap or synthetic-answer style could drive the vector and metric.  A semantic preference proxy is not sufficient evidence of clinical safety or factual correctness.

\textbf{Suggested fix.} Document the full benchmark pipeline and create a group split by source passage and drug/entity.  Release split IDs and overlap checks.  Add blinded clinician adjudication and structured error scoring, report tie counts and complete BERTScore configuration, and use ``RefPref rate'' consistently for the current endpoint.

\subsection*{9. Schedule evidence and generation-cap effects remain conflated}

\textbf{Classification: claim/evidence mismatch.}

\textbf{Location.} Abstract; both primary evaluations; Table~VII; conclusion.

\textbf{Problem.} In natural completion, Hard Cutoff is the only schedule with a Holm-adjusted significant gain; Linear Decay gains 2.40 points with $p_{\text{adj}}=0.1189$.  The bounded test instead highlights Linear Decay, but its EOS column is 0.0\%, indicating that all 500 outputs are truncated at 200 tokens.  The baseline itself changes from 73.00\% at 200 tokens to 65.40\% under natural completion, showing that the endpoint is sensitive to the cap.

The umbrella phrase ``early-stopping steering'' alternates between Linear Decay alone and both Linear Decay and Hard Cutoff.  The factorial table contains very small metric differences but no paired uncertainty for schedule comparisons.

\textbf{Why it matters.} The bounded Linear Decay result cannot be assumed to transfer to completed answers, and the natural experiment supports Hard Cutoff rather than the specifically promoted decay schedule.

\textbf{Suggested fix.} Pre-specify a primary schedule and endpoint.  Call the 200-token analysis a truncated-prefix stress test, report output-length distributions, and provide paired schedule comparisons.  State plainly that the current natural-completion evidence favors Hard Cutoff and treats Linear Decay as exploratory.

\subsection*{10. Statistical interval descriptions and layer-wise conclusions still need revision}

\textbf{Classification: likely statistical-reporting issue and unsupported inference.}

\textbf{Location.} Table~III caption; contributions; RAG comparison; Layer-Wise Subspace Probing.

\textbf{Problem.} The exact McNemar $p$-values shown for the baseline comparisons are reproducible from the discordant counts.  However, Table~III calls its intervals a paired $B=10{,}000$ percentile bootstrap with seed 42, while straightforward paired percentile resampling of the displayed four-outcome counts does not reproduce the quoted endpoints; several endpoints resemble analytic approximation intervals.  The code and per-question data are not provided to resolve the discrepancy.  The Wilcoxon test concerns continuous BERTScore values, not the binary RefPref endpoint, and should not be presented as a second confirmation of the same quantity.

Separately, a 0.7040--0.7140 validation range across layers on $N_{\text{eval}}=100$, without paired uncertainty, does not establish that truthfulness is distributed across layers or ``confirm framework robustness.''

\textbf{Why it matters.} The statistical workflow is a named contribution, and the layer claim turns an exploratory selection sweep into an unsupported representation-level conclusion.

\textbf{Suggested fix.} Release per-example outputs and exact analysis code; name the interval method precisely; state Wilcoxon zero handling and the tested BERTScore quantity; and define the family covered by Holm adjustment.  For the layer sweep, report all paired results and uncertainty and limit the conclusion to the observation that several tested layers had similar validation scores.

\section*{Compilation, editorial, and presentation issues}

\begin{enumerate}
 \item \textbf{The current source does not compile cleanly.} Line 123 contains eleven literal \verb|\n| control sequences.  A \texttt{pdflatex -halt-on-error} build stops at the first one with ``Undefined control sequence''.  Replace them with actual source line or paragraph breaks.
 \item \textbf{Missing math delimiters.} The phrases \verb|(+5.15\sigma)| and \verb|(+2.52\sigma / +2.39~pp)| use \verb|\sigma| outside math mode, producing two ``Missing \$ inserted'' errors and corrupting the following paragraph's typesetting.  Enclose each expression in math delimiters.
 \item \textbf{Duplicated setup text.} The full environment/software/decoding sentence appears at the end of line 123 and is repeated immediately on line 124.  Keep one version after resolving the dtype contradiction.
 \item \textbf{Duplicated discussion.} The unpenalized Pareto-trade-off discussion appears twice in the same paragraph.  Keep one concise version and separate any verified repetition-penalty experiment into Results.
 \item \textbf{Unreferenced tables.} The natural-completion contingency table and the new Hybrid RAG contingency table have labels but are not explicitly cited in the surrounding prose.  Cite and interpret both tables.
 \item \textbf{Table readability.} Table~V compresses seven metric columns into one IEEE column, and the Hybrid contingency table also overflows its column.  Use \verb|table*|, split metrics, or shorten headings rather than relying on extreme scaling.
\end{enumerate}

\section*{Bounded proofing sweep}

\begin{itemize}
 \item \textbf{RAG paragraph:} ``achieving Top-1 Recall ... and MRR ..., achieving 76.20\% RefPref'' repeats ``achieving''.  Replace the second instance with ``and yields''.
 \item \textbf{Latency units:} write $7.32$~s and $+7.13$~s consistently rather than \texttt{7.32s} and ``+7.13s''.
 \item \textbf{Natural-completion repetition:} 30.5\% is the relative increase for Hard Cutoff; Linear Decay increases from 4.10\% to 5.37\%, approximately 31.0\%.  Name the schedule when giving the percentage.
 \item \textbf{Reference list:} every cited key has a matching \verb|\bibitem|, and no obvious duplicated punctuation or malformed quotation punctuation was found.  Standardize edition style, for example ``2nd ed.'' if IEEE conventions are followed.
\end{itemize}

The mandatory pattern scan returned two candidates, both false positives after inspection: the mathematical ellipsis in $\{1,\ldots,100\}$ and the metric name ``Recall@3'' after a semicolon.  No inverse-trigonometric or coordinate-transform formula appears, so no branch or quadrant ambiguity was identified.

\section*{Highest-priority fixes}

\begin{enumerate}
 \item Make the source compile without ignored errors, then reconcile dtype, layer index, and autoregressive step semantics.
 \item Replace the contradictory activation-norm descriptions with one documented experiment and narrow the mechanistic conclusion.
 \item Reconcile the covariance-control run count and replace the invariant Pair-Shuffled control.
 \item Correct or remove the undocumented repetition-penalty result and its $11\times$ arithmetic error.
 \item Report that steering and Hybrid RAG are statistically indistinguishable on RefPref, and frame their result as a systems trade-off.
 \item Strengthen source-disjoint evaluation and clinician adjudication before making factual or clinical claims.
\end{enumerate}

\section*{Items requiring external verification}

The following points cannot be established from the manuscript alone:

\begin{itemize}
 \item the actual run configuration and logs needed to determine whether \texttt{float16} or \texttt{bfloat16} was used, which decoder block was hooked, and how generation steps were counted;
 \item whether 20 or 100 covariance-matched vectors were executed and whether the retained summary values come from the expanded run;
 \item the existence and complete results of the claimed repetition-penalty ``Experiment 10'';
 \item expert involvement, factual validation, public availability, and redistribution rights for ViHaluEval-Medical and the formulary-derived passages;
 \item whether \texttt{ref3} and \texttt{ref11} support the specific clinical-error and internal-activation claims attributed to them.
\end{itemize}

\section*{Internally consistent checks}

The dataset split counts sum to 14,700, the three evaluation categories sum to 500, the category-level preferred counts reproduce the stated totals, and the dose formulas
\[
 D_{\text{decay}}=\frac{K+1}{2}|\alpha_0|,
 \qquad
 \alpha_{\text{match}}=0.0425\alpha_0
 \quad(K=16,T=200)
\]
are correct.  The displayed baseline-comparison McNemar $p$-values agree with their discordant counts, and the new Hybrid table's margins correctly reproduce 386/500 and 381/500.  No dimensional, sign, inverse-trigonometric, or branch error was found in the schedule equations themselves.  These correct checks do not resolve the reporting and implementation contradictions above.

\end{document}